\section{Exact Gaussian models and calibration geometry}
\label{sec:gaussian-geometry}

Gaussian reservoirs and Gaussian descriptions of homogeneous coexistence
peaks have a substantial history~\cite{ChallaHetherington1988PRL,ChallaLandauBinder1986}.
Here they are exact probability models used to isolate calibration and
phase-selection effects. Their interphase tails are not substituted for
the microscopic laws of Section~\ref{sec:microscopic}. A quadratic
residual weight also differs from the exact power-law reservoir: the
identification $\kappa=\beta^2/(C_B/k_B)$ describes a local entropy curvature,
not a globally constant heat capacity for an exactly quadratic entropy.

\subsection{Two scalar phases and the optimized error crossover}

Let $\varphi_{s^2}$ denote a centered normal density with standard
deviation $s>0$, and consider
\begin{equation}
 p_N(E)=\tfrac12\varphi_{s_N^2}(E-m_N)+\tfrac12\varphi_{s_N^2}(E+m_N),
 \quad m_N/s_N\to\infty,
 \qquad
 q_{N,t}\propto p_N e^{tE-\kappa_NE^2/2},\quad\kappa_N\ge0.
 \label{eq:scalar-gaussian-model}
\end{equation}
The field $t\in\R$ is unrestricted. Completion of squares gives the
transformed component variance, means, and upper-phase weight:
\begin{equation}
 \widetilde s_N^2=\frac{s_N^2}{1+\kappa_Ns_N^2},\qquad
 \widetilde m_{\pm,N}=\frac{\pm m_N+t s_N^2}{1+\kappa_Ns_N^2},\qquad
 v_{+,N}(t)=\frac1{1+e^{-2tm_N/(1+\kappa_Ns_N^2)}}.
 \label{eq:scalar-gaussian-completion}
\end{equation}

\begin{proposition}[Scalar convergence and optimized crossover]
\label{prop:scalar-gaussian}
There exists a field sequence giving $\TV(p_N,q_{N,t_N})\to0$ if and
only if $\kappa_Nm_Ns_N\to0$. If instead
$\kappa_Nm_Ns_N\to a\in[0,\infty)$, then
\begin{align}
 \TV(p_N,q_{N,0})&\longrightarrow2\Phi(a/2)-1,
 \label{eq:scalar-balanced-crossover}\\
 \inf_{t\in\R}\TV(p_N,q_{N,t})
       &\longrightarrow\min\{2\Phi(a/2)-1,1/2\}.
 \label{eq:scalar-optimal-crossover}
\end{align}
\end{proposition}
\begin{proof}
Put $r_N=m_N/s_N$, $k_N=\kappa_Ns_N^2$, and $b_N=t_Ns_N$.
In units of $s_N$, the transformed means are
$(\pm r_N+b_N)/(1+k_N)$ and their variance is $(1+k_N)^{-1}$.
If $\kappa_Nm_Ns_N\to0$, then $k_N\to0$ and the choice $t_N=0$
preserves the weights exactly while making both standardized mean
displacements vanish. Componentwise Gaussian TV convergence proves
sufficiency.

For necessity, we use the scalar case of the label-matching argument
given in full for Theorem~\ref{thm:gaussian-geometry}. Fix $R>0$ and
take the windows $|E\mp m_N|\le Rs_N$. Each has target mass tending to
$\tfrac12[2\Phi(R)-1]$, and the windows are disjoint for large $N$.
A transformed component has standard deviation
$\widetilde s_N\le s_N$, so the smaller of its probabilities in
these two windows tends to zero as $m_N/s_N\to\infty$.
TV convergence therefore forces both
components to keep weights with lower limit at least
$\tfrac12[2\Phi(R)-1]$, each supplying one window, and
letting $R\to\infty$ gives weights tending to $1/2$ and centers
$\widetilde m_{\pm,N}=\pm m_N+O(s_N)$, matched in their natural
order. Subtracting these center estimates and using the exact
transformed means gives
\[
 \frac{2r_N}{1+k_N}=2r_N+O(1),\qquad
 \frac{k_Nr_N}{1+k_N}=O(1).
\]
Since $r_N\to\infty$, this forces $k_N\to0$, so the variance ratio
already tends to one. Restricting to the matched windows and then
letting $R\to\infty$ gives TV convergence of each transformed
component to its target normal. In particular, its probability above
the target mean tends to $1/2$. With standardized displacement
$\delta_{\pm,N}=(b_N\mp k_Nr_N)/(1+k_N)$, this probability is
$\Phi(\delta_{\pm,N}\sqrt{1+k_N})$. Hence both displacements tend
to zero. Subtracting them proves $k_Nr_N\to0$.

For the finite crossover, $k_N\to0$ follows already from $k_Nr_N\to a$
and $r_N\to\infty$. At $t_N=0$ the two translated limiting normals
have displacements $\mp a$ and unit variance. Separating their windows
gives~\eqref{eq:scalar-balanced-crossover}, using
$\TV(\Normal(0,1),\Normal(a,1))=2\Phi(|a|/2)-1$.

To optimize, take any sequence of fields and a subsequence on which
the transformed upper weight converges. If its limit lies strictly
between zero and one, the logistic formula in
\eqref{eq:scalar-gaussian-completion} implies $b_N=O(1/r_N)$.
The limiting component displacements are therefore still $\mp a$.
The limiting TV is a convex function of the limiting upper weight,
and reflection makes that function symmetric about $1/2$. It is
minimized there, giving the first term in~\eqref{eq:scalar-optimal-crossover}.
If one weight tends to zero, a single bounded-width Gaussian cannot
cover both target windows, so the limiting lower bound is $1/2$.
This also covers means that escape all target windows. These subsequence
bounds apply to arbitrary approximate minimizers of the infimum.
The balanced choice attains the first bound. For $a>0$, the choice
$t_N=\kappa_Nm_N$ centers the upper component exactly at $m_N$, makes
its weight tend to one, and attains $1/2$. If $a=0$, the balanced bound
already vanishes.
\end{proof}

The optimizer changes behavior at
$a_*=2\Phi^{-1}(3/4)\simeq1.34898$. Above this value the TV objective
can favor an accurate single phase while discarding the other one.
An optimized scalar distance should therefore be accompanied by the
phase probabilities.

Temperature calibration changes the required scale even within this
simple model. Give the upper reference phase weight $w\ne1/2$, with
$0<w<1$, and hold the tangent of the quadratic bath at the canonical
mean $\bar E_N=(2w-1)m_N$. This means
$q_N\propto p_Ne^{-\kappa_N(E-\bar E_N)^2/2}$, or
$t_N=\kappa_N\bar E_N$. The exact transformed phase odds are
\begin{equation}
 \log\frac{v_{+,N}}{1-v_{+,N}}
 =\log\frac w{1-w}
   +\frac{2\kappa_N(2w-1)m_N^2}{1+\kappa_Ns_N^2}.
 \label{eq:gaussian-untuned-odds}
\end{equation}
The same separated-window argument first requires
$\kappa_Ns_N^2\to0$; convergence of the weights then requires
$\kappa_Nm_N^2\to0$. Conversely, that condition makes both the weight
and shape errors vanish. Thus for $m_N\asymp N$, $s_N\asymp\sqrt N$,
this specified untuned prescription requires $\kappa_NN^2\to0$,
whereas an appropriately tuned two-phase law needs only
$\kappa_NN^{3/2}\to0$. The reference point in this comparison is fixed
by the target canonical mean, not by a self-consistent modified mean.

\subsection{Several exchanged quantities and phase geometry}

Let the Gaussian variable be $Y\in\R^d$, where here $d$ is the number
of exchanged quantities, not a spatial dimension. Fix distinct phase
points $e_1,\ldots,e_k$ and positive weights $w_i$ summing to one. Set
\begin{equation}
 p_N=\sum_{i=1}^k w_i\Normal(Ne_i,NI_d),\qquad
 q_{N,t}\propto p_N\exp(t\cdot Y-\kappa_N|Y|^2/2),
 \quad t\in\R^d,\quad\kappa_N\ge0.
 \label{eq:multivariate-gaussian-model}
\end{equation}
With
\begin{equation}
 s_N=(1+\kappa_NN)^{-1},\qquad
 \alpha_N=\frac{\kappa_NN^2}{1+\kappa_NN},\qquad
 z_N=\frac{Nt}{1+\kappa_NN},
 \label{eq:multivariate-parameters}
\end{equation}
the transformed covariance is $Ns_NI_d$, its means are $Ns_N(e_i+t)$,
and its weights are
\begin{equation}
 v_{i,N}=\frac{w_i\exp(z_N\cdot e_i-\alpha_N|e_i|^2/2)}
                 {\sum_jw_j\exp(z_N\cdot e_j-\alpha_N|e_j|^2/2)}.
 \label{eq:multivariate-weights}
\end{equation}
For any affine dependence $(a_1,\ldots,a_k)$, meaning
$\sum_i a_i=0$ and $\sum_i a_ie_i=0$, this gives the exact invariant
\begin{equation}
 \sum_i a_i\log(v_{i,N}/w_i)
       =-\frac{\alpha_N}{2}\sum_i a_i|e_i|^2.
 \label{eq:multivariate-invariant}
\end{equation}
Every adjustable field has canceled.

\begin{theorem}[Geometry of Gaussian reservoir convergence]
\label{thm:gaussian-geometry}
For $k\ge2$, some field sequence in
\eqref{eq:multivariate-gaussian-model} gives TV convergence if and only if
\begin{equation}
 \begin{cases}
 \kappa_NN^{3/2}\to0,&\text{if the phase points are cospherical},\\
 \kappa_NN^2\to0,&\text{otherwise}.
 \end{cases}
 \label{eq:gaussian-geometry-classification}
\end{equation}
For $k=1$, the condition is $\kappa_NN\to0$. Cospherical means
that the points lie on a common sphere of positive radius in $\R^d$;
in $d=1$ a sphere is a pair of points, so three distinct scalar
phase points are never cospherical.
\end{theorem}
\begin{proof}
First justify matching the component labels. On the macroscopic
variable $Y/N$, every transformed covariance is at most $I_d/N$.
If TV tends to zero, each target phase neighborhood of positive mass
must therefore receive a distinct transformed component of nonvanishing
weight. There are exactly $k$ components, so this is a bijection, and
the matched transformed centers $s_N(e_i+t_N)$ lie in bounded
neighborhoods of the phase points; in particular they are bounded.
Along any subsequence, take a
further subsequence on which $s_N$, the translation $s_Nt_N$, and the
matching permutation converge. The limiting map $e\mapsto se+a$
maps the finite phase set onto itself. Its positive diameter forces
$s=1$, and the arithmetic centroid then forces $a=0$. Distinct phase
points make the matching permutation the identity. It follows that
$s_N\to1$, $s_Nt_N\to0$, and $v_{i,N}\to w_i$ with the original labels.

Isolating each matched component in expanding fluctuation windows
then gives its Gaussian TV convergence. In particular,
\begin{equation}
 \sqrt N\{(s_N-1)e_i+s_Nt_N\}\longrightarrow0
 \quad\hbox{for every }i.
 \label{eq:gaussian-shape-necessity}
\end{equation}
This necessity can also be detected by projection onto the direction
of a component's mean displacement. Subtracting two equations in
\eqref{eq:gaussian-shape-necessity} proves
$\kappa_NN^{3/2}\to0$ for any two distinct phase points.

The phase-weight equations $v_i=w_i$, when $\alpha_N>0$, are solvable
exactly if and only if
\begin{equation}
 |e_i|^2=2c\cdot e_i+b\quad\hbox{for every }i
 \label{eq:cosphere-equation}
\end{equation}
for some $c\in\R^d,b\in\R$. This is the sphere condition. Taking
$t_N=\kappa_NNc$ then restores every weight exactly and makes the
standardized mean changes
$\kappa_NN^{3/2}(c-e_i)/(1+\kappa_NN)$ tend to zero. This proves
sufficiency in the spherical case.

If~\eqref{eq:cosphere-equation} has no solution, finite-dimensional
linear algebra gives an affine dependence with
$\sum_i a_i|e_i|^2\ne0$. Since $v_{i,N}/w_i\to1$,
\eqref{eq:multivariate-invariant} forces $\alpha_N\to0$, equivalent
to $\kappa_NN^2\to0$. Taking $t_N=0$ proves sufficiency, including
the component shapes. For a single phase, $t_N=\kappa_NNe_1$ aligns
its mean exactly; matching the covariance is equivalent to
$\kappa_NN\to0$ and is necessary whatever mean is chosen.
\end{proof}

Every affinely independent set of at most $d+1$ points is cospherical,
but larger sets may also be cospherical. The number of phases alone
therefore does not determine the threshold in this vector model.
For scalar phase points $-1,0,1$, the invariant reads
\begin{equation}
 \frac{v_{+,N}v_{-,N}}{v_{0,N}^2}
       =\frac{w_+w_-}{w_0^2}e^{-\alpha_N}.
 \label{eq:three-gaussian-invariant}
\end{equation}
Three distinct scalar points are not cospherical, so their threshold
is $\kappa_NN^2\to0$.

\subsection{The exact optimum when some phases must disappear}

An empty-sphere contact set is a nonempty subset of phase indices of
the form
\begin{equation}
 F(c)=\operatorname*{arg\,min}_{1\le i\le k}|e_i-c|^2,
 \qquad c\in\R^d.
 \label{eq:empty-sphere-contact}
\end{equation}
These are the contact sets of the usual lower-paraboloid, or Delaunay,
construction \cite{EdelsbrunnerSeidel1986}. Concentration of exponential
families on faces also has an established general theory
\cite{CsiszarMatus2005}; the theorem here determines the exact optimized
TV loss under the specified diverging quadratic parameter.
Only finitely many subsets can occur. Define
$W_*:=\max_c\sum_{i\in F(c)}w_i$.

\begin{theorem}[Optimal Gaussian phase loss]
\label{thm:gaussian-phase-loss}
For the model~\eqref{eq:multivariate-gaussian-model}, if
\begin{equation}
 \kappa_NN^{3/2}\to0,\qquad \kappa_NN^2\to\infty,
 \label{eq:gaussian-phase-loss-regime}
\end{equation}
then
\begin{equation}
 \lim_{N\to\infty}\inf_{t\in\R^d}\TV(p_N,q_{N,t})=1-W_*.
 \label{eq:gaussian-phase-loss-optimum}
\end{equation}
The optimization includes field sequences of unbounded norm.
\end{theorem}
\begin{proof}
For a fixed center $c$, choose $t_N=\kappa_NNc$. The phase score in
\eqref{eq:multivariate-weights} is, up to a common constant,
$-\alpha_N|e_i-c|^2/2$. Since $\alpha_N\to\infty$, the weights
converge to the reference weights conditioned on $F(c)$. The first
condition in~\eqref{eq:gaussian-phase-loss-regime} makes every
component shape change vanish in TV. This attains the upper bound
$1-\sum_{F(c)}w_i$, and maximizing over contact sets proves the desired
upper bound.

For the lower bound, take any sequence of fields. First consider a
subsequence with $|t_N|\ge\varepsilon>0$. Put
$u_N=t_N/|t_N|$ and $H_N=\max_i u_N\cdot e_i$.
The weight outside indices with
$u_N\cdot e_i\ge H_N-N^{-1/4}$ tends to zero: relative to a maximizer
its linear score loses at least $c_1N^{3/4}$ for a fixed $c_1>0$,
whereas its quadratic
score differs by only $O(\alpha_N)=o(\sqrt N)$. For the remaining
indices the projected transformed mean per particle is at least
\[
 s_N(H_N-N^{-1/4}+|t_N|)\ge H_N+\varepsilon/2
\]
for all large $N$, uniformly also when $|t_N|\to\infty$.
The target Gaussian mixture has negligible mass beyond
$u_N\cdot Y/N>H_N+\varepsilon/4$, while the transformed law has mass
tending to one there. Its TV therefore tends to one.

It remains to consider $t_N\to0$. Pass to a subsequence on which all
weights converge, and let $I$ be the nonempty set of indices with
positive limiting weights. The transformed means divided by $N$
approach their own $e_i$, since $s_N\to1$. Fixed disjoint small
macroscopic neighborhoods of the phase points then imply
\begin{equation}
 \liminf\TV(p_N,q_{N,t_N})\ge1-\sum_{i\in I}w_i.
 \label{eq:phase-loss-missing-mass}
\end{equation}
To constrain $I$, fix $i_0\in I$ and put $c_N=z_N/\alpha_N$.
The exact score ratios imply
\begin{align*}
 c_N\cdot(e_i-e_{i_0})
       -\tfrac12(|e_i|^2-|e_{i_0}|^2)&\to0&& (i\in I),\\
 c_N\cdot(e_j-e_{i_0})
       -\tfrac12(|e_j|^2-|e_{i_0}|^2)&\le o(1)&& (1\le j\le k).
\end{align*}
For the inequalities use $v_{i_0,N}$ bounded below and $v_{j,N}\le1$;
the resulting log-ratio upper bound divided by $\alpha_N$ vanishes.
These are approximate equalities and inequalities of a fixed finite
linear system. Such a system cannot be infeasible while its right-hand
side errors tend to zero. Indeed, Farkas' alternative would supply fixed
multipliers annihilating the unknown vector and yielding a strictly
negative constant, contradicting the vanishing errors. Thus there is
a finite vector $c$ satisfying the limiting system exactly, even if
$c_N$ has no bounded subsequence. Its sphere contact set $F(c)$ contains
$I$. Hence $\sum_{i\in I}w_i\le W_*$. Combining with
\eqref{eq:phase-loss-missing-mass} proves the lower bound along every
subsequence of approximate minimizers, and completes the proof.
\end{proof}

For three scalar points $-1,0,1$, the largest possible retained sets
are the two adjacent pairs. Thus the limiting optimal error is
$\min(w_-,w_+)$, including $1/3$ for equal weights. The central phase
cannot be the only discarded phase. More generally,
Theorem~\ref{thm:gaussian-phase-loss} identifies the exact retained-mass
criterion, including the field sequences whose sphere centers escape
to infinity; bounded-center analysis alone would supply only an upper
bound.

\subsection{Anisotropic curvature: the effective metric}

For the same common covariance $NI_d$, replace $\kappa_NI_d$ by a
symmetric positive-semidefinite matrix $B_N$. Completing squares gives
\begin{equation}
 C_N=N(I_d+NB_N)^{-1},\qquad
 M_N=B_N(I_d+NB_N)^{-1},\qquad
 z_N=N(I_d+NB_N)^{-1}t_N,
 \label{eq:anisotropic-effective-metric}
\end{equation}
as the transformed covariance, effective metric, and effective field.
The means are $N(I_d+NB_N)^{-1}(e_i+t_N)$, and
$v_i\propto w_i\exp[z_N\cdot e_i-N^2e_i^{\mathsf T}M_Ne_i/2]$.
Accordingly the affine invariants and exact weight-restoration equations
use $M_N$, not the raw matrix $B_N$. Completing the square for
positive-definite $M_N$ gives an ellipsoid, with a single point in the
zero-radius case. For positive-rank singular $M_N$, a linear coefficient
in its range gives an elliptic cylinder at positive radius or an affine
subspace at zero radius. A nonzero linear component in its nullspace
instead gives a paraboloid, possibly with additional flat directions.
For $M_N=0$, a nonzero linear coefficient gives a hyperplane; if the
coefficient and compatible constant both vanish, the locus is the whole
space. Inconsistent level constants give an empty locus in the
completed-square or constant cases. These are finite-$N$ algebraic statements.
The scalar exponent classification above is not asserted for arbitrary
anisotropic sequences or unequal phase covariance matrices.
